
\rev{\textbf{Experimental Setup.}} Clinical note mining for rare disease differential diagnosis involves two key tasks: extracting phenotypes and identifying rare diseases. \rev{Here, we summarize the experimental setup towards evaluating existing mining approaches and RDMA.}

\textbf{Datasets.} \rev{Existing rare disease datasets have known limitations but remain useful benchmarks. Phenotype identification. We evaluate on the case study cohort (CSC) \cite{garcia2024_HPORAG} and BioLark gold standard corpora \cite{yang2023enhancingphenotyperecognitionclinical}. The CSC benchmark better reflects clinical notes with more text and phenotypes per document, though it still falls short of real-world clinical complexity. Rare disease identification. Three benchmarks exist: RDD \cite{fabregat2018deep_rdd}, RareDis \cite{martinez2022raredis}, and a noisily annotated MIMIC-III disease benchmark \cite{johnson2016mimic3,dong2023ontology_poor_annotations}. We exclude RDD as it lacks rare disease mention annotations. To demonstrate RDMA's ability to both improve extraction and assist annotators, we construct a gold evaluation set from the MIMIC-III annotations via combined human and agent annotation, yielding a realistic benchmark on real clinical notes. All benchmark statistics are described in Supplementary Table~\sTabDatasets.}

\rev{\textbf{Evaluation.} Our goal is to map text evidence to structured codes from the HPO \cite{gargano2024mode_hpo} and Orphanet \cite{weinreich2008orphanet} ontologies, framing extraction as a medical coding task \cite{Edin_2023_medical_coding_review}. Phenotype extraction. Annotated codes are available for both phenotype datasets, making evaluation straightforward. Rare disease identification. RareDis provides only text span annotations without ontology mappings. As both exist in our MIMIC3 evaluation set, we report both span and code mapping performance here as MIMIC3-RD Entity and MIMIC3-RD Code respectively. For baselines lacking an ontology mapping step, we choose to use a fuzzy matching retrieval approach that is employed by the authors of PhenoGPT \cite{yang2023enhancingphenotyperecognitionclinical}.} All evaluation metrics are described in the Supplementary Material.


\textbf{Baselines.} \rev{Across all benchmarks, approaches fall into three categories: (1) rule-based methods that tokenize text and match to ontologies via nearest-neighbor search, (2) fine-tuned encoder models trained to extract entities directly, and (3) LLM-based methods using zero-shot or RAG-based extract-and-match pipelines.
For rule-based approaches, we include a simple retrieve-and-string-match Dictionary Match baseline using state-of-the-art clinical embeddings MedEmbed \cite{balachandran2024medembed} and the state-of-the-art FastHPOCR. For fine-tuned encoders, we distinguish between off-the-shelf models and those we fine-tune ourselves. Off-the-shelf, we evaluate PhenoGPT \cite{yang2023enhancingphenotyperecognitionclinical} and Stanza's i2b2 Clinical BERT, both used without modification. Following the RareDis authors \cite{segura2022exploring_raredis}, we additionally fine-tune BioBERT \cite{sun2021biomedical_biobert_ner} and BioClinicalBERT \cite{alsentzer2019publiclyavailableclinicalbert} for rare disease and phenotype extraction; all four encoder models are evaluated across all benchmark datasets. For LLM-based approaches, we test a range of models in both zero-shot and RAG-based settings \cite{garcia2024_HPORAG}. We note that highly specialized approaches such as FastHPOCR and PhenoGPT do not apply to our rare disease extraction benchmarks.}